\documentclass[10pt,a4paper]{amsart}
\usepackage[margin=19mm]{geometry}
\usepackage{amsmath,amssymb,mathtools}
\usepackage[scr=rsfs,cal=euler]{mathalfa}
\usepackage{xfrac}
\usepackage[unicode,hidelinks,bookmarks=false]{hyperref}
\newcommand{\ie}{i.e.\ }
\newcommand{\cf}{cf.\ }
\newcommand{\resp}{respectively\ }
\newcommand{\Subsection}{\subsection}
\providecommand{\nobreakdash}{\nobreak\hbox{-}}
\setlength{\emergencystretch}{2em}
\multlinegap0pt
\usepackage{bm}
\usepackage{bbm}
\usepackage{dsfont}
\usepackage{xspace}
\usepackage{cancel}
\usepackage{mathtools}
\usepackage{amsthm}
\makeatletter
\@ifundefined{thm}{\newtheorem{thm}{Thm}}{}
\@ifundefined{definition}{\newtheorem{definition}[thm]{Definition}}{}
\@ifundefined{lemma}{\newtheorem{lemma}[thm]{Lemma}}{}
\makeatother
\newcommand{\R}{{\mathbb R}}
\newcommand{\p}{{ \partial}}
\title[The semilinear heat inequality with Morrey initial data on R]{The semilinear heat inequality with Morrey initial data on Riemannian manifolds}
\author{Anuk Dayaprema}
\date{Source: arXiv:2412.21029v2 (CC BY 4.0)}
\begin{document}
\maketitle

\noindent\emph{Source and licence.} The semilinear heat inequality with Morrey initial data on Riemannian manifolds. Version arXiv:2412.21029v2 (CC BY 4.0), distributed under the Creative Commons Attribution 4.0 International licence, as recorded in the arXiv licence metadata for this exact version and shown on the abstract page https://arxiv.org/abs/2412.21029v2. Full rights notice: licenses/arxiv-2412.21029-CC-BY-4.0.txt.\par\smallskip

\noindent\textit{Editorial scope.} Counted unit: the Lemma whose source label is \texttt{lemma:checkinghmfisMorreysolution} in the article (source0.tex lines 1359--1364, source label \texttt{lemma:checkinghmfisMorreysolution}), delivered together with the article's own complete original proof of it (source0.tex lines 1365--1481, one \texttt{proof} environment of 117 lines). No mathematics is written, repaired or re-derived: statement and proof are the article's own text line for line. Required context retained uncounted: diffineq (lines 299--301); defn:weaksolution (lines 357--376); lemma:gradientestimates (lines 1347--1355). Every definition, display and label of those lines is retained unchanged. Omitted from this package and not redistributed: 1--298, 302--356, 377--1346, 1356--1358, 1482--1777. Nothing inside a retained excerpt is omitted or altered: every retained body is delivered exactly as the article's lines are written, so no omission receipt of any kind accompanies this package. Citations: the retained excerpts carry no citation command at all, so no bibliography entry of the article is transcribed and no reference is redistributed. Reference adaptation: none was needed and none was made; every reference inside the delivered unit resolves inside it, so no unresolved reference or citation is produced. Printed slips kept literal: no printed word, symbol, hypothesis or number of the counted statement or of its proof was altered. Layout and class: the article's own class is \texttt{amsart} with packages including amssymb, amsmath, amscd, amsfonts, latexsym, verbatim, mathrsfs, color, array, xy, multirow, graphicx, amsthm, bbm; the wrapper is the lane's pinned \texttt{amsart} editorial wrapper at 10pt on A4, which restates the theorem-like environments the retained text needs and adds the article's own macro definitions that the retained text needs (\texttt{\textbackslash R}, \texttt{\textbackslash p}), with extra wrapper packages (bm, bbm, dsfont, xspace, cancel, mathtools). The delivered numbering is the unit's own. Assets: no retained span contains a figure environment, a graphics-inclusion command, a PDF-page inclusion, a table or a TikZ picture; no image file is copied into this package, the package declares no asset dependency and no image file is redistributed with it. Licence: version arXiv:2412.21029v2 of this article is distributed under the Creative Commons Attribution 4.0 International licence, as recorded in the arXiv licence metadata for this exact version and shown on the abstract page https://arxiv.org/abs/2412.21029v2. A full rights notice is delivered at licenses/arxiv-2412.21029-CC-BY-4.0.txt. Characters: every retained excerpt is pure ASCII, so the delivered engine, XeLaTeX with its default Latin Modern text font, reports no missing character.
\begin{align}\label{diffineq}
    \left( \frac{\p}{\p t} - \Delta \right) u \leq Au^p + Bu,
\end{align} with nonnegative constants \(A\) and \(B\). These flows often exist smoothly as long as the energy density remains bounded; hence, 

\begin{definition}\label{defn:weaksolution}
        Let \(p, q, r\), and \(T\) be constants such that
        \begin{align*}
            p > 1 + \frac{2}{n}, \quad 1 < q \leq q_c := \frac{n(p-1)}{2}, \quad 1 \leq \frac{r}{p} < q < r, \quad T > 0.
        \end{align*}Set
        \begin{align*}
            \lambda := \frac{2q}{p-1}.
        \end{align*} We say \(u\) is a solution to (\ref{diffineq}) on \([0, T)\) if:
        \begin{enumerate}
            \item \(u \in C_{\text{loc}}^0([0, T), M^{q, \lambda}) \cap C_{\text{loc}}^0((0, T), L^\infty)\). 
            \item \(u(t)\) satisfies 
        \begin{align*}
            \lim_{t \searrow 0} \left(\|u(t) - u(0)\|_{M^{q, \lambda}} + t^{\frac{1}{p-1}} \|u(t)\|_\infty + t^{\frac{\lambda}{2}\left(\frac{1}{q} - \frac{1}{r}\right)}\|u(t)\|_{M^{r, \lambda}}\right)  = 0.
        \end{align*}
        \item For each \(t \in [0, T)\), we have \(u(t) \geq 0\) a.e., and for \(0 \leq t' < t < T\) we have a.e.
        \begin{align}\label{weaksolution:subsolution}
            u(t) \leq H_{t-t'}(u(t')) + \int_{t'}^t H_{t-s}(Au^p(s) + Bu(s)) \, ds.
        \end{align}
        \end{enumerate}
    \end{definition}

   \begin{lemma}\label{lemma:gradientestimates}
        Let \(0 < t \leq 1\) and \(1 \leq q \leq r \leq \infty\). Then, there exists \(C \geq 1\), depending only on the geometry of \(M\), such that
   \begin{align*}
       %&\|\nabla e^{t\Delta}f_0\|_r \leq C t^{-\frac{n}{2}\left(\frac{1}{q} -\frac{1}{r}\right) - \frac{1}{2}} \|f_0\|_q, \\
       &\|\nabla H_t(f_0)\|_r \leq C t^{-\frac{n}{2}\left(\frac{1}{q} -\frac{1}{r}\right)} \cdot \begin{cases} t^{-\frac12}\|f_0\|_q & f_0 \in L^q\\
       \| \nabla f_0\|_q & f_0 \in W^{1,q}.
       \end{cases}
   \end{align*}
   \end{lemma}

\begin{lemma}\label{lemma:checkinghmfisMorreysolution}
    We have \(w \in C^0([0, T], W^{1, n}) \cap C_{\text{loc}}^0((0, T], W^{1, \infty})\). Moreover, 
    \begin{align*}
        \lim_{t \searrow 0} \left(\|w(t) - w_0\|_{W^{1, n}} + t^{\frac{1}{6}}\|\nabla w(t)\|_{\frac{3n}{2}} + t^{\frac{1}{2}}\|\nabla w(t)\|_{\infty}\right) = 0.
    \end{align*} In particular, \(|\nabla w|\) is a solution of (\ref{diffineq}) with \(p = 3, q = 2\), and \(r = 3\) as in Def. \ref{defn:weaksolution}.
\end{lemma}

\begin{proof}
   
   Since \(w(t)\) is smooth for \(t > 0\) and since \(M\) is closed, we have
   \begin{align*}
       w \in C_{\text{loc}}^0((0, T], W^{1, \infty}) \subset C_{\text{loc}}^0((0, T], W^{1, n}).
   \end{align*} Thus, we just show that the limit in the statement of this lemma is zero. We establish this in the following order:
   \begin{enumerate}
       \item \(\lim_{t \searrow 0} t^{\frac{1}{6}} \|\nabla w(t)\|_{\frac{3n}{2}} = 0\); 
       \item \(\lim_{t \searrow 0} t^{\frac{1}{2}} \|\nabla w(t)\|_\infty = 0\); 
       \item \(\lim_{t \searrow 0} \|w(t) - w_0\|_{W^{1, n}} = 0\).
   \end{enumerate}\par
   In the sequel, \(C\) denotes a positive constant depending on the geometries of \(M\) and \(N\) which may increase from line to line. \par
   \emph{Item 1: \(\lim_{t \searrow 0} t^{\frac{1}{6}} \|\nabla w(t)\|_{\frac{3n}{2}} = 0\).}
   Denote
   \begin{align*}
       w_1(t) := H_t(w_0).
   \end{align*} We first claim that
   \begin{align}\label{checkinghmfisMorreysolution:goestozero}
      \lim_{t \searrow 0} \left(t^{\frac{1}{6}} \|\nabla w_1(t)\|_{\frac{3n}{2}} + t^{\frac{1}{2}}\|\nabla w_1(t)\|_\infty\right) = 0. 
   \end{align} To see this, first note that since \(w_0 \in W^{1, n}\) and since \(M\) is closed, there exists a sequence \(\tilde{w}_i \in C^\infty(M, \R^k)\) converging to \(w_0\) in \(W^{1, n}(M, \R^k)\). Then by Lemma \ref{lemma:gradientestimates},
    \begin{align*}
        \|\nabla w_1(t)\|_{\frac{3n}{2}} &\leq \|\nabla w_1(t) - \nabla H_t(\tilde{w}_i)\|_{\frac{3n}{2}} + \|\nabla H_t(\tilde{w}_i)\|_{\frac{3n}{2}} \\
        &\leq C t^{-\frac{1}{6}}\|\nabla w_0 - \nabla \tilde{w}_i\|_n + C \|\nabla H_t(\tilde{w}_i)\|_\infty \\
        &\leq C t^{-\frac{1}{6}}\|\nabla w_0 - \nabla \tilde{w}_i\|_n + C \|\nabla \tilde{w}_i\|_\infty.
    \end{align*} We multiply by \(t^{\frac{1}{6}}\). Choosing \(i\) and then \(t\) appropriately, we deduce that \(\lim_{t \searrow 0} t^{\frac{1}{6}} \|\nabla w_1(t)\|_{\frac{3n}{2}} = 0\). An analogous estimation implies that \(\lim_{t \searrow 0} t^{\frac{1}{2}} \|\nabla w_1(t)\|_{\infty} = 0\). Thus, the claim (\ref{checkinghmfisMorreysolution:goestozero}) holds. \par
    We now use (\ref{checkinghmfisMorreysolution:goestozero}) to prove Item 1. By Wang's construction, \(w = \lim_{n \to \infty} \Phi(w_n)\) in \(X_T\), where \(w_1\) is as before and \(w_n = \Phi(w_{n-1})\).
    Set
    \begin{align}\label{checkinghmfisMorreysolution:duhamelconstant}
        C_{(\ref{checkinghmfisMorreysolution:duhamelconstant})} := \int_0^1 (1-\sigma)^{-\frac{5}{6}} \sigma^{-\frac{2}{3}} \, d\sigma < \infty.
    \end{align} 
    Recall that we have assumed \(T \leq 1\). Then for all \(i\) and each \(\tau \in (0, T]\), Lemma \ref{lemma:gradientestimates} yields
    \begin{align*}
        \|\nabla w_{i+1}(\tau)\|_{\frac{3n}{2}} - \|\nabla w_1(\tau)\|_{\frac{3n}{2}} &\leq  \int_0^\tau \|\nabla H_{\tau-s}( \Pi(w_i(s))(\nabla w_i(s), \nabla w_i(s)))\|_{\frac{3n}{2}} \, ds \\ 
        &\leq C\int_0^\tau (\tau-s)^{-\frac{5}{6}}\|\Pi(w_i(s))(\nabla w_i(s), \nabla w_i(s))\|_{\frac{3n}{4}} \, ds \\ 
        &\leq  C\int_0^\tau (\tau-s)^{-\frac{5}{6}}\|\nabla w_i(s)\|_{\frac{3n}{2}}^2 \, ds \\ 
        &\leq  C\left(\sup_{0 < s \leq \tau} s^{\frac{1}{6}} \|\nabla w_i(s)\|_{\frac{3n}{2}}\right)^2\int_0^\tau (\tau-s)^{-\frac{5}{6}}s^{-\frac{1}{3}} \, ds \\ 
        &\leq  C\left(\sup_{0 < s \leq \tau} s^{\frac{1}{6}} \|\nabla w_i(s)\|_{\frac{3n}{2}}\right)^2C_{(\ref{checkinghmfisMorreysolution:duhamelconstant})} \tau^{-\frac{1}{6}}.
    \end{align*} Hence,
    \begin{align*}
        \tau^{\frac{1}{6}} \|\nabla w_{i+1}(\tau)\|_{\frac{3n}{2}} \leq \sup_{0 < s \leq \tau} s^{\frac{1}{6}}\|\nabla w_1(s)\|_{\frac{3n}{2}} + CC_{(\ref{checkinghmfisMorreysolution:duhamelconstant})}\left(\sup_{0 < s \leq \tau} s^{\frac{1}{6}} \|\nabla w_i(s)\|_{\frac{3n}{2}}\right)^2.
    \end{align*} Fix \(t \in (0, T]\). The preceding inequality holds for all \(\tau \in (0, t]\), so we replace \(\sup_{0 < s \leq \tau}\) in the RHS by \(\sup_{0 < s \leq t}\) and then supremize the LHS over \(\tau \in (0, t]\) to obtain%. We deduce
    \begin{align}\label{checkinghmfisMorreysolution:item1inductioninequality}
        \sup_{0 < s \leq t}s^{\frac{1}{6}} \|\nabla w_{i+1}(s)\|_{\frac{3n}{2}} \leq \sup_{0 < s \leq t} s^{\frac{1}{6}}\|\nabla w_1(s)\|_{\frac{3n}{2}} + CC_{(\ref{checkinghmfisMorreysolution:duhamelconstant})}\left(\sup_{0 < s \leq t} s^{\frac{1}{6}} \|\nabla w_i(s)\|_{\frac{3n}{2}}\right)^2.
    \end{align} We now use (\ref{checkinghmfisMorreysolution:item1inductioninequality}) and induction to show that \(\sup_{0 < s \leq t} s^{\frac{1}{6}} \|\nabla w_i(s)\|_{\frac{3n}{2}}\) can be made small for all \(i\) if \(t\) is chosen small enough. \par 
    Let
    \begin{align*}
        \varepsilon \in \left(0, (4CC_{(\ref{checkinghmfisMorreysolution:duhamelconstant})})^{-1}\right).
    \end{align*}(\ref{checkinghmfisMorreysolution:goestozero}) implies there exists \(t > 0\) such that 
    \begin{align*}
        \sup_{0 < s \leq t} s^{\frac{1}{6}}\|\nabla w_1(s)\|_{\frac{3n}{2}} \leq \varepsilon.
    \end{align*} Denote
    \begin{align*}
        \mu_i := \sup_{0 < s \leq t} s^{\frac{1}{6}} \|\nabla w_i(s)\|_{\frac{3n}{2}}.
    \end{align*} By choice of \(t\), we have
    \begin{align*}
        \mu_1 \leq \varepsilon.
    \end{align*} Suppose that we have 
    \begin{align*}
        \mu_i \leq 2\varepsilon.
    \end{align*} Then by (\ref{checkinghmfisMorreysolution:item1inductioninequality})
    \begin{align*}
        \mu_{i+1} \leq \varepsilon + CC_{(\ref{checkinghmfisMorreysolution:duhamelconstant})}\mu_i^2 \leq \varepsilon + 4CC_{(\ref{checkinghmfisMorreysolution:duhamelconstant})}\varepsilon^2 \leq 2\varepsilon.
    \end{align*} Thus by induction, \(\mu_i \leq 2\varepsilon\) for all \(i\).  From the definition of the \(X_T\) norm and the convergence \(w_i \to w\) in \(X_T\), we have that \(\nabla w_i \to \nabla w\) in \(C_{\text{loc}}^0((0, T), L^{\infty})\). Hence, 
    \begin{align*}
        \sup_{0 < s \leq t} s^{\frac{1}{6}}\|\nabla w(s)\|_{\frac{3n}{2}}\leq 2\varepsilon,
    \end{align*} so Item 1 follows. \par
    \emph{Item 2: \(\lim_{t \searrow 0} t^{\frac{1}{2}}\|\nabla w(t)\|_\infty = 0\).}  Since \(w\) solves the fixed-point equation, we have by Lemma \ref{lemma:gradientestimates}
    \begin{align*}
        \|\nabla w(t)\|_\infty - \|\nabla w_1(t)\|_\infty  &= \|\nabla \Phi(w)(t)\|_\infty - \|\nabla w_1(t)\|_\infty \\
        &\leq C\int_0^t (t-s)^{-\frac{5}{6}} \|\Pi(w(s))(\nabla w(s), \nabla w(s)) \|_{\frac{3n}{2}} \, ds \\
        &\leq C\left(\sup_{0 < s \leq t} s^{\frac{1}{2}} \|\nabla w(s)\|_\infty\right)\left(\sup_{0 < s \leq t} s^{\frac{1}{6}} \|\nabla w(s)\|_{\frac{3n}{2}}\right)\int_0^t (t-s)^{-\frac{5}{6}} s^{-\frac{2}{3}} \, ds \\
        &\leq C\left(\sup_{0 < s \leq t} s^{\frac{1}{2}} \|\nabla w(s)\|_\infty\right)\left(\sup_{0 < s \leq t} s^{\frac{1}{6}} \|\nabla w(s)\|_{\frac{3n}{2}}\right)C_{(\ref{checkinghmfisMorreysolution:duhamelconstant})}t^{-\frac{1}{2}}. 
    \end{align*} We multiply by \(t^{\frac{1}{2}}\). By definition of the \(X_T\) norm, we know that \(\sup_{0 < s \leq T} s^{\frac{1}{2}}\|\nabla w(s)\|_\infty < \infty\). Then, since by Item 1 and (\ref{checkinghmfisMorreysolution:goestozero})
    \begin{align*}
        \lim_{t \searrow 0} \left(t^{\frac{1}{6}}\|\nabla w(t)\|_{\frac{3n}{2}} + t^{\frac{1}{2}}\|\nabla w_1(t)\|_\infty\right) = 0,
    \end{align*} Item 2 follows. \par
    \emph{Item 3: \(\lim_{t \searrow 0} \|w(t) - w_0\|_{W^{1, n}} = 0\).} Since \(w_0 \in W^{1, n}\), \(w_1(t) \to w_0\) in \(W^{1, n}(M, \R^k)\) as \(t \searrow 0\). Next, we estimate
    \begin{align*}
        \left\|\int_0^t H_{t-s}(\Pi(w)(\nabla w(s), \nabla w(s))) \, ds\right\|_n &\leq C\int_0^t (t-s)^{-\frac{1}{3}} \|\Pi(w)(\nabla w(s), \nabla w(s))\|_{\frac{3n}{4}} \, ds \\
        &\leq C \left(\sup_{0 < s \leq t} s^{\frac{1}{6}} \|\nabla w(s)\|_{\frac{3n}{2}}\right)^2 \int_0^t (t-s)^{-\frac{1}{3}}s^{-\frac{1}{3}} \, ds \\
        &\leq CC_{(\ref{checkinghmfisMorreysolution:duhamelconstant})} \left(\sup_{0 < s \leq t} s^{\frac{1}{6}} \|\nabla w(s)\|_{\frac{3n}{2}}\right)^2
    \end{align*} and
    \begin{align*}
        \left\|\int_0^t \nabla H_{t-s}(\Pi(w)(\nabla w(s), \nabla w(s))) \, ds\right\|_n &\leq \int_0^t C(t-s)^{-\frac{2}{3}} \|\Pi(w)(\nabla w(s), \nabla w(s))\|_{\frac{3n}{4}} \, ds \\
        &\leq CC_{(\ref{checkinghmfisMorreysolution:duhamelconstant})}  \left(\sup_{0 < s \leq t} s^{\frac{1}{6}} \|\nabla w(s)\|_{\frac{3n}{2}}\right)^2.
    \end{align*} Since \(w\) solves the fixed-point equation, we deduce that 
    \begin{align*}
        \|w(t) - w_0\|_{W^{1, n}} &= \|\Phi(w)(t) - w_0\|_{W^{1, n}} \\
        &\leq \|w_1(t) - w_0\|_{W^{1, n}} + \left\|\int_0^t  H_{t-s}(\Pi(w(s))(\nabla w(s), \nabla w(s))) \,ds \right\|_{W^{1, n}} \\
        &\to 0 \ \text{ as } \ t \searrow 0,
    \end{align*} as desired.
    \par
    We have thus shown that the limit in the statement of the lemma is zero. It remains to deduce that
    \begin{align*}
        u := |\nabla w|
    \end{align*}is a solution of (\ref{diffineq}), which we now explain. It follows from H{\"o}lder's inequality that
    \begin{align*}
        u \in C^0([0, T], M^{2, 2}) \cap  C_{\text{loc}}^0((0, T], L^\infty),
    \end{align*} and furthermore
    \begin{align*}
       \lim_{t \searrow 0} \left(t^{\frac{1}{6}}\|u\|_{M^{3, 2}} + t^{\frac{1}{2}}\|u\|_\infty\right) = 0. 
    \end{align*} Since \(w\) is smooth on \(M \times (0, T]\), the B{\"o}chner identity for harmonic map flow and the comparison principle yields for all \(0 < t' < t \leq T\) 
    \begin{align}\label{checkinghmfisMorreysolution:subsolution}
        u(t) \leq H_{t-t'}(u(t')) + \int_{t'}^t H_{t-s}(Au^3(s) + Bu(s)) \, ds.
    \end{align} Note that for fixed \(t > 0\), we have that as \(t' \searrow 0\)
    \begin{align*}
        \| H_{t-t'}(u(t')) - H_t(u(0))\|_n &= \|H_{t-t'}(u(t')-H_{t'}(u(0)))\|_n \\
        &\leq C\|u(t') - H_{t'}(u(0))\|_n \\
        &\leq C(\|u(t') - u(0)\|_n + \|u(0) - H_{t'}(u(0))\|_n) \\
        &\to 0
    \end{align*} and 
    \begin{align*}
        \left\|\int_0^{t'} \hspace{-0.4em} H_{t-s} (Au^3(s) + Bu(s)) \, ds\right\|_n &\leq CA\hspace{-0.4em}\int_0^{t'} \hspace{-0.4em} (t-s)^{-1}\|u^3(s)\|_{\frac{n}{3}} \, ds + CB\hspace{-0.4em}\int_0^{t'} \hspace{-0.4em}\|u(s)\|_n \, ds \\
        &\leq CA(\|u(0)\|_n+1)^3\ln\left(\frac{t}{t-t'}\right) + CB(\|u(0)\|_n+1)t' \\
        &\to 0.
    \end{align*} Therefore, (\ref{checkinghmfisMorreysolution:subsolution}) also holds with \(t' = 0\). Hence, we conclude that \(|\nabla w|\) is a solution of (\ref{diffineq}).
\end{proof}

\end{document}
